% ============================================================================
%  ESTADÍSTICA INFERENCIAL — Sesión 3
%  Semana 3: 24 al 29 de agosto de 2026
%  1.3 Regla de Laplace y enfoques para asignar probabilidades
%  1.4 Axiomas de la probabilidad y reglas de la probabilidad
% ============================================================================
\documentclass[aspectratio=169,10pt]{beamer}

\usepackage[utf8]{inputenc}
\IfFileExists{spanish.ldf}{\usepackage[spanish,es-noshorthands]{babel}}{}
\usepackage{amsmath,amssymb}
\usepackage{booktabs}
\usepackage{tikz}

\usetheme{Uniminuto}

\uniasignatura{Estadística inferencial}
\uniprofesor{Andrés Felipe Puerta Velez}
\unicorreo{andres.puerta-v@uniminuto.edu.co}
\uninrc{CBASBL061}
\unisemana{Sesión 3 — Semana 3 · 24 al 29 de agosto de 2026}

\newcommand{\listatiny}{\setbeamerfont{itemize/enumerate body}{size=\tiny}\tiny}

\begin{document}

\portada

\begin{frame}{Orden del día}
  \begin{enumerate}
    \item Llamado a lista y revisión de los ejercicios propuestos
    \item Repaso: dónde quedamos
    \item \azul{Tema 1.3: enfoques para asignar probabilidades y regla de Laplace}
    \item \azul{Tema 1.4: axiomas y reglas de la probabilidad}
    \item Ejercicios en clase
    \item Cierre y ejercicios propuestos
  \end{enumerate}
\end{frame}

\begin{frame}{Repaso: dónde quedamos}
  Ya sabemos \azul{describir} eventos (sesión 1) y \azul{contarlos} (sesión 2).
  \[ P(A)=\frac{\#A}{\#\Omega} \]

  \vspace{0.2cm}
  \begin{alertblock}{Las preguntas de hoy}
    \footnotesize
    ¿De dónde sale esa fórmula y cuándo tenemos derecho a usarla? ¿Qué pasa si
    los resultados \emph{no} son igualmente probables? ¿Qué reglas debe cumplir
    cualquier asignación de probabilidad para ser coherente?
  \end{alertblock}
\end{frame}

\seccion{1.3 y 1.4 Probabilidad: enfoques, axiomas y reglas}

% ------------------------------------------------------------- ENFOQUES ----
\begin{frame}{Tres enfoques para asignar probabilidades}
  \footnotesize
  \begin{center}
  \begin{tabular}{@{}p{2.4cm}p{4.6cm}p{4.6cm}@{}}
    \toprule
    \textbf{Enfoque} & \textbf{Cómo se asigna} & \textbf{Cuándo se usa} \\
    \midrule
    Clásico \newline (a priori) &
      $P(A)=\#A/\#\Omega$, suponiendo resultados igualmente probables &
      Juegos de azar, sorteos, muestreo aleatorio: la simetría está garantizada
      por el diseño \\
    \addlinespace
    Frecuentista \newline (a posteriori) &
      $P(A)\approx f_r(A)=\dfrac{\text{veces que ocurrió }A}{\text{número de repeticiones}}$ &
      Cuando hay datos históricos: tasas de defectos, de morosidad, de rotación \\
    \addlinespace
    Subjetivo &
      Grado de creencia de un experto, coherente con los axiomas &
      Eventos que no se repiten: lanzar un producto nuevo, ganar una licitación \\
    \bottomrule
  \end{tabular}
  \end{center}

  \vspace{0.15cm}
  {\scriptsize\color{uniGris} Los tres producen números entre 0 y 1 y obedecen
  las mismas reglas; lo que cambia es de dónde sale el número\,[2].}
\end{frame}

\begin{frame}{Regla de Laplace}
  \begin{block}{Enunciado}
    \small
    Si $\Omega$ es finito y todos sus resultados son \destacar{igualmente
    probables}, entonces para cualquier evento $A$
    \[ P(A)=\frac{\#A}{\#\Omega}
          =\frac{\text{casos favorables}}{\text{casos posibles}} \]
  \end{block}

  \vspace{0.2cm}
  \begin{alertblock}{La condición no es un detalle}
    \footnotesize
    Si los resultados no son equiprobables, la fórmula da respuestas
    equivocadas. Al lanzar dos dados y sumar, $\Omega=\{2,3,\dots,12\}$ tiene 11
    resultados, pero $P(\text{suma}=7)\neq 1/11$: hay que trabajar con los 36
    pares ordenados, que sí son equiprobables.
  \end{alertblock}
\end{frame}

\begin{frame}{Ejercicio 1 — en clase}
  \begin{alertblock}{Individual (6 minutos)}
    Se extrae una carta al azar de una baraja inglesa de \destacar{52 cartas}
    (13 de cada palo: corazones, diamantes, tréboles y picas).

    \vspace{0.3cm}
    \begin{enumerate}
      \item $P(\text{sale un as})$
      \item $P(\text{sale un corazón})$
      \item $P(\text{sale un as \emph{o} un corazón})$
      \item ¿Por qué el punto 3 no es simplemente la suma de 1 y 2?
    \end{enumerate}
  \end{alertblock}
\end{frame}

\begin{frame}{Ejercicio 1 — solución}
  \begin{columns}[T,onlytextwidth]
    \begin{column}{0.48\textwidth}
      \begin{block}{Cálculos}
        \small
        \begin{align*}
          P(\text{As}) &= \tfrac{4}{52}=\tfrac{1}{13}=0{,}0769\\[0.25em]
          P(\heartsuit) &= \tfrac{13}{52}=0{,}25\\[0.25em]
          P(\text{As}\cup\heartsuit) &= \tfrac{4}{52}+\tfrac{13}{52}-\tfrac{1}{52}
             =\tfrac{16}{52}=0{,}3077
        \end{align*}
      \end{block}
    \end{column}
    \begin{column}{0.48\textwidth}
      \begin{exampleblock}{Punto 4}
        \footnotesize
        Porque el \azul{as de corazones} pertenece a los dos eventos. Si se
        sumaran sin más, esa carta se contaría dos veces.

        \vspace{0.2cm}
        Es la misma corrección de la sesión 1:
        \[ \#(A\cup B)=\#A+\#B-\#(A\cap B) \]
      \end{exampleblock}
    \end{column}
  \end{columns}
\end{frame}

% ------------------------------------------------------- FRECUENTISTA -----
\begin{frame}{Enfoque frecuentista}
  Cuando no hay simetría, la probabilidad se estima observando:
  \[ P(A)\;\approx\; f_r(A)=\frac{n_A}{n} \]
  donde $n_A$ es el número de veces que ocurrió $A$ en $n$ repeticiones.

  \vspace{0.2cm}
  \begin{exampleblock}{Ejemplo}
    \footnotesize
    De 2\,000 pedidos despachados el semestre pasado, 160 llegaron tarde:
    $P(\text{tarde})\approx 160/2000 = 0{,}08$.
  \end{exampleblock}

  \vspace{0.15cm}
  \begin{alertblock}{Ley de los grandes números}
    \footnotesize
    A medida que $n$ crece, la frecuencia relativa $f_r(A)$ se estabiliza
    alrededor de la probabilidad verdadera. Es la justificación de fondo de
    todo el curso: \destacar{con muestras grandes, lo observado se parece a lo
    real}.
  \end{alertblock}
\end{frame}

% --------------------------------------------------------------- AXIOMAS ---
\begin{frame}{Axiomas de la probabilidad (Kolmogórov)}
  \begin{block}{Toda asignación de probabilidad debe cumplir}
    \small
    \begin{enumerate}
      \item \textbf{No negatividad:} $P(A)\ge 0$ para todo evento $A$.
      \item \textbf{Normalización:} $P(\Omega)=1$.
      \item \textbf{Aditividad:} si $A$ y $B$ son mutuamente excluyentes
            ($A\cap B=\varnothing$), entonces $P(A\cup B)=P(A)+P(B)$.
    \end{enumerate}
  \end{block}

  \vspace{0.25cm}
  {\footnotesize\color{uniGris} Son solo tres, y de ellos se deduce todo lo
  demás. Cualquier número que no los respete —una probabilidad negativa, o dos
  eventos excluyentes cuyas probabilidades sumen más de 1— indica un error de
  cálculo\,[1].}
\end{frame}

\begin{frame}{Consecuencias de los axiomas}
  \small
  \begin{columns}[T,onlytextwidth]
    \begin{column}{0.48\textwidth}
      \begin{block}{Reglas que se deducen}
        \begin{align*}
          P(\varnothing) &= 0\\[0.2em]
          P(A^{c}) &= 1-P(A)\\[0.2em]
          0 \le P(A) &\le 1\\[0.2em]
          \text{si } A\subseteq B:\; P(A) &\le P(B)
        \end{align*}
      \end{block}
    \end{column}
    \begin{column}{0.48\textwidth}
      \begin{exampleblock}{La regla del complemento es la más útil}
        \footnotesize
        Calcular ``al menos uno'' directamente obliga a sumar muchos casos.
        Calcular ``ninguno'' suele ser un solo caso.

        \vspace{0.2cm}
        \[ P(\text{al menos un 6})=1-P(\text{ningún 6}) \]
        \[ = 1-\frac{25}{36}=\frac{11}{36}=0{,}3056 \]
      \end{exampleblock}
    \end{column}
  \end{columns}
\end{frame}

\begin{frame}{Regla de la suma}
  \begin{block}{Caso general}
    \[ P(A\cup B)=P(A)+P(B)-P(A\cap B) \]
  \end{block}

  \begin{block}{Si $A$ y $B$ son mutuamente excluyentes}
    \[ P(A\cup B)=P(A)+P(B) \qquad\text{porque } P(A\cap B)=0 \]
  \end{block}

  \vspace{0.2cm}
  \begin{alertblock}{Para tres eventos}
    \footnotesize
    \[ P(A\cup B\cup C)=P(A)+P(B)+P(C)-P(A\cap B)-P(A\cap C)-P(B\cap C)
       +P(A\cap B\cap C) \]
    Se suman los individuales, se restan los pares y se vuelve a sumar el
    triple: es el principio de inclusión--exclusión.
  \end{alertblock}
\end{frame}

\begin{frame}{Ejercicio 2 — en clase}
  \begin{alertblock}{En parejas (7 minutos)}
    Una empresa tiene \destacar{500 empleados}: 300 hombres y 200 mujeres. De
    ellos, \destacar{150 tienen posgrado} (90 hombres y 60 mujeres). Se elige
    un empleado al azar.

    \vspace{0.25cm}
    \begin{enumerate}
      \item $P(H)$, que sea hombre
      \item $P(G)$, que tenga posgrado
      \item $P(H\cap G)$
      \item $P(H\cup G)$
      \item $P(G^{c})$, que no tenga posgrado
    \end{enumerate}
  \end{alertblock}
\end{frame}

\begin{frame}{Ejercicio 2 — solución}
  \begin{columns}[T,onlytextwidth]
    \begin{column}{0.44\textwidth}
      \footnotesize
      \begin{center}
      \begin{tabular}{@{}lccc@{}}
        \toprule
         & Posgrado & No & Total \\
        \midrule
        Hombres & 90  & 210 & 300 \\
        Mujeres & 60  & 140 & 200 \\
        \midrule
        Total   & 150 & 350 & 500 \\
        \bottomrule
      \end{tabular}
      \end{center}
    \end{column}
    \begin{column}{0.52\textwidth}
      \begin{block}{Cálculos}
        \small
        \begin{align*}
          P(H) &= 300/500 = \azul{0{,}60}\\
          P(G) &= 150/500 = \azul{0{,}30}\\
          P(H\cap G) &= 90/500 = \azul{0{,}18}\\
          P(H\cup G) &= 0{,}60+0{,}30-0{,}18 = \azul{0{,}72}\\
          P(G^{c}) &= 1-0{,}30 = \azul{0{,}70}
        \end{align*}
      \end{block}
    \end{column}
  \end{columns}

  \vspace{0.1cm}
  {\footnotesize\color{uniGris} Organizar los datos en una \azul{tabla de doble
  entrada} evita casi todos los errores. La usaremos de nuevo en la sesión 4.}
\end{frame}

\begin{frame}{Ejercicio 3 — en clase}
  \begin{alertblock}{Individual (8 minutos)}
    Se lanzan \destacar{dos dados} corrientes y se observa el par de
    resultados. Recuerde que $\#\Omega=36$.

    \vspace{0.3cm}
    \begin{enumerate}
      \item $P(\text{la suma sea }7)$
      \item $P(\text{aparezca al menos un }6)$
      \item $P(\text{la suma sea al menos }10)$
      \item Resuelva el punto 2 usando la regla del complemento.
    \end{enumerate}
  \end{alertblock}
\end{frame}

\begin{frame}{Ejercicio 3 — solución}
  \begin{columns}[T,onlytextwidth]
    \begin{column}{0.52\textwidth}
      \begin{block}{Conteos sobre los 36 pares}
        \footnotesize
        \begin{align*}
          \text{Suma }7&:\;\{(1,6),(2,5),(3,4),(4,3),(5,2),(6,1)\}\\
          &\Rightarrow P=\tfrac{6}{36}=\azul{0{,}1667}\\[0.3em]
          \text{Al menos un }6&:\; 11 \text{ pares}
            \Rightarrow P=\tfrac{11}{36}=\azul{0{,}3056}\\[0.3em]
          \text{Suma}\ge 10&:\; 3+2+1=6
            \Rightarrow P=\tfrac{6}{36}=\azul{0{,}1667}
        \end{align*}
      \end{block}
    \end{column}
    \begin{column}{0.44\textwidth}
      \begin{exampleblock}{Punto 4}
        \footnotesize
        ``Ningún 6'' significa que ambos dados caen en $\{1,\dots,5\}$:
        $5\times 5=25$ pares.
        \[ P=1-\frac{25}{36}=\frac{11}{36} \]
        Mucho más rápido que enumerar los 11 casos favorables.
      \end{exampleblock}
    \end{column}
  \end{columns}
\end{frame}

\begin{frame}{Resumen de reglas}
  \footnotesize
  \begin{center}
  \begin{tabular}{@{}lll@{}}
    \toprule
    \textbf{Regla} & \textbf{Expresión} & \textbf{Cuándo} \\
    \midrule
    Laplace          & $P(A)=\#A/\#\Omega$ & resultados equiprobables \\
    Frecuentista     & $P(A)\approx n_A/n$ & hay datos históricos \\
    Rango            & $0\le P(A)\le 1$ & siempre \\
    Complemento      & $P(A^{c})=1-P(A)$ & siempre; ideal para ``al menos uno'' \\
    Suma general     & $P(A\cup B)=P(A)+P(B)-P(A\cap B)$ & siempre \\
    Suma excluyentes & $P(A\cup B)=P(A)+P(B)$ & si $A\cap B=\varnothing$ \\
    \bottomrule
  \end{tabular}
  \end{center}

  \vspace{0.25cm}
  \begin{alertblock}{Lo que todavía no sabemos hacer}
    \footnotesize
    Calcular $P(A\cap B)$ cuando no viene dado. Para eso necesitamos la
    \destacar{probabilidad condicional}, tema de la sesión 4.
  \end{alertblock}
\end{frame}

\begin{frame}{Cierre de la sesión}
  \begin{block}{Lo que debe quedar claro hoy}
    \begin{itemize}
      \item Hay tres formas de asignar probabilidades: clásica, frecuentista y
            subjetiva; las tres obedecen los mismos axiomas.
      \item Laplace solo aplica si los resultados son equiprobables.
      \item $P(A^{c})=1-P(A)$ resuelve casi todos los ``al menos uno''.
      \item En la suma siempre se descuenta la intersección, salvo que los
            eventos sean excluyentes.
    \end{itemize}
  \end{block}

  \vspace{0.2cm}
  \begin{alertblock}{Trabajo independiente para la próxima sesión}
    \footnotesize
    Resolver los \destacar{20 ejercicios de la lámina siguiente}; las
    respuestas están en la lámina posterior. Leer el capítulo de axiomas y
    reglas de la probabilidad de la referencia [1].
  \end{alertblock}
\end{frame}

\begin{frame}{Ejercicios propuestos}
  \vspace{-0.1cm}
  \begin{beamercolorbox}[wd=\textwidth,sep=0.14cm]{block body}
    \centering\scriptsize
    \textbf{Del 1 al 10:} encuesta a 400 clientes de un banco. 240 usan la
    aplicación ($A$), 180 tienen tarjeta de crédito ($T$) y 120 tienen ambas.
  \end{beamercolorbox}
  \vspace{0.1cm}
  \listatiny
  \vspace{0.15cm}
  \begin{columns}[T,onlytextwidth]
    \begin{column}{0.47\textwidth}
      \begin{enumerate}\setlength{\itemsep}{0.16cm}
        \item $P(A)$
        \item $P(T)$
        \item $P(A\cap T)$
        \item $P(A\cup T)$
        \item $P(A^{c})$
        \item $P(T^{c})$
        \item $P\big((A\cup T)^{c}\big)$
        \item $P(A\cap T^{c})$
        \item $P(T\cap A^{c})$
        \item $P(A\,\triangle\,T)$
      \end{enumerate}
    \end{column}
    \begin{column}{0.50\textwidth}
      \begin{enumerate}\setlength{\itemsep}{0.16cm}
        \setcounter{enumi}{10}
        \item Un dado: $P(\text{par})$
        \item Un dado: $P(\text{mayor que }4)$
        \item Un dado: $P(\text{par o mayor que }4)$
        \item Dos dados: $P(\text{suma}=7)$
        \item Dos dados: $P(\text{suma}\ge 10)$
        \item Dos dados: $P(\text{al menos un }6)$
        \item Baraja: $P(\text{figura: J, Q o K})$
        \item Baraja: $P(\text{as o rey})$
        \item Moneda 3 veces: $P(\text{exactamente 2 caras})$
        \item Moneda 3 veces: $P(\text{al menos una cara})$
      \end{enumerate}
    \end{column}
  \end{columns}
\end{frame}

\begin{frame}{Respuestas}{Autoevaluación de los ejercicios propuestos}
  \listatiny
  \vspace{0.15cm}
  \begin{columns}[T,onlytextwidth]
    \begin{column}{0.47\textwidth}
      \begin{enumerate}\setlength{\itemsep}{0.16cm}
        \item $0{,}60$
        \item $0{,}45$
        \item $0{,}30$
        \item $0{,}75$
        \item $0{,}40$
        \item $0{,}55$
        \item $0{,}25$
        \item $0{,}30$
        \item $0{,}15$
        \item $0{,}45$
      \end{enumerate}
    \end{column}
    \begin{column}{0.50\textwidth}
      \begin{enumerate}\setlength{\itemsep}{0.16cm}
        \setcounter{enumi}{10}
        \item $3/6=0{,}5000$
        \item $2/6=0{,}3333$
        \item $4/6=0{,}6667$ \;{\color{uniGris}(se resta el 6)}
        \item $6/36=0{,}1667$
        \item $6/36=0{,}1667$
        \item $11/36=0{,}3056$
        \item $12/52=0{,}2308$
        \item $8/52=0{,}1538$
        \item $3/8=0{,}3750$
        \item $7/8=0{,}8750$ \;{\color{uniTitulo}$=1-1/8$}
      \end{enumerate}
    \end{column}
  \end{columns}
\end{frame}

\begin{frame}{Referencias}
  \footnotesize
  \begin{itemize}
    \item[{[1]}] Díaz Rodríguez, M. (2022). \textit{Estadística inferencial
      aplicada}. Universidad del Norte.
      \texttt{https://elibro.net/es/lc/uniminuto/titulos/221614}
    \item[{[2]}] Llinás Solano, H. (2014). \textit{Introducción a la
      estadística con aplicaciones en ciencias sociales}. Universidad del
      Norte. \texttt{https://elibro.net/es/lc/uniminuto/titulos/70062}
    \item[{[3]}] Contento Rubio, M. R. (2019). \textit{Estadística con
      aplicaciones en R}. Editorial Utadeo.
      \texttt{https://elibro.net/es/lc/uniminuto/titulos/220926}
  \end{itemize}

  \vspace{0.35cm}
  \begin{center}
    {\Large\color{uniAzul}\textbf{¿Preguntas?}}
  \end{center}
\end{frame}

\end{document}
